\documentclass[11pt]{article}
\usepackage{amsmath,amssymb,amsthm,graphicx}
\newtheorem{theorem}{Theorem}
\newcommand{\Z}{\mathbb{Z}}
\title{Random walks on a fictitious lattice}
\author{Ada Example\\ University of Nowhere \and Bo Sample\\ Institute of Placeholders}
\date{Version 2, 15 September 2024}
\begin{document}
\maketitle
\begin{abstract}
We study the simple random walk $X_n$ on $\Z^2$ and prove that it returns to
the origin infinitely often. This paper is an \emph{example} of an OpenSheaf
deposit: its authors and results are invented.
\end{abstract}
\section{Introduction}
Let $(X_n)_{n\ge 0}$ be the simple random walk on $\Z^2$ started at $0$,
as in Figure~\ref{fig:walk}; see \cite{Polya,Sample} for imaginary background.
\begin{figure}[h]\centering\includegraphics[width=0.45\textwidth]{figures/walk.pdf}
\caption{A path of the walk.}\label{fig:walk}\end{figure}
\begin{theorem}\label{thm:main}
$\mathbb{P}(X_n = 0 \text{ for infinitely many } n) = 1.$
\end{theorem}
\begin{proof}
Since $\mathbb{P}(X_{2n}=0)\sim \frac{1}{\pi n}$, the expected number of
returns $\sum_n \mathbb{P}(X_{2n}=0)$ diverges, and the walk is recurrent.
\end{proof}
\section{Remarks}The constant $\frac{1}{\pi n}$ follows from Stirling's formula. We thank an imaginary referee for this remark.
\bibliographystyle{plain}
\bibliography{refs}
\end{document}
